\section{The $V$-Hierarchy}

\begin{boxdefinition}[The $V$-Hierarchy]
    Define $V_0 = \emptyset$, $V_{\alpha + 1} = \Powset\of{V_{\alpha}}$ for any ordinal $\alpha$, and $V_{\beta} = \bigcup_{\alpha < \beta} V_{\alpha}$ for any limit ordinal $\beta$.
\end{boxdefinition}

One can show that $\alpha \le \beta \to V_{\alpha} \subseteq V_{\beta}$.

We can ask ourselves what exactly ZFC is describing. The answer is, \textbf{the $V$-hierarchy}.

Interestingly,
\begin{align*}
    V_{\omega} \models \ZFC - \text{infinity} + \neg \text{infinity}
\end{align*}

Even more interestingly, ZFC doesn't prove its own consistency (if it could, it would be inconsistent, by Gödel), but it comes exactly one axiom away: it can prove the existence of models of everything except for the axiom of infinity.



\subsection{The Existence of $V_{\omega}$}

It is worth asking ourselves the following question: \textit{Why does $V_{\omega}$ exist?} Or, more precisely, \textit{why is $V_{\omega}$ a set?}

By induction, we know that each $V_n$ is a set. We want to essentially see that the map $n \mapsto V_n$ is a \textit{set} (ie, a \textit{function}). Then, if we take the union of the range of this function, whose domain is the set $\omega$ whose existence is given by \Cref{ZFC:Infty}, we can define $V_{\omega}$ to be that union.

We invoke \Cref{ZFC:Replacement}, the Axiom of Replacement. Let $\varphi\of{x, y}$ be the formula capturing the following facts:
\begin{enumerate}
    \item $y$ is a function with domain $\dom{y} = x \cup \set{x}$.
    \item If $\emptyset \in \dom{y}$, then $y\of{\emptyset} = \emptyset$.
    % [CORRECTED] was: $y\of{z \cup \set{z}} = \Powset\of{z}$ --- the recursion has to take the power set of the previous \textit{value}, $y\of{z} = V_z$, not of the index $z$; with $\Powset\of{z}$ the sequence agrees with the $V_n$ only up to $n = 3$ (where $n = V_n$) and gives $8$ elements rather than $16$ at $n = 4$
    \item If $z \cup \set{z} \in \dom{y}$, then $z \in \dom{y}$ and $y\of{z \cup \set{z}} = \Powset\of{y\of{z}}$.
\end{enumerate}
The idea is that $y = \cycl{V_0, V_1, V_2, \ldots, V_{n}}$\footnote{By this, we mean it is a \textit{sequence}, ie, a function from $\omega$ to something that takes $0$ to $V_{0}$, $1$ to $V_{1}$, and so on.} for some $n \in \omega$. Then,
\begin{align*}
    \ZFC \vdash \forall n \in \omega \ \exists! y \, \varphi\of{n, y}
\end{align*}
Applying the Axiom of Replacement (\Cref{ZFC:Replacement}), we know that the map that takes any $n \in \omega$ to its unique witness $y$ for $n$ \textit{is, indeed, a function}. In other words,
\begin{align*}
    \grpres{\grpres{V_m}{m \leq n}}{n < \omega}
\end{align*}
is a set. From this, we can read off
\begin{align*}
    \grpres{V_n}{n < \omega}
\end{align*}
by taking a union. % OF WHAT?

\begin{remark}
    Recall that $\parenth{V_{\omega}, \in}$ is shorthand for $\parenth{V_{\omega}, R}$, where $R = \setst{\parenth{x, y} \in V_{\omega} \times V_{\omega}}{x \in y}$. It is possible to show not only that
    \begin{align*}
        \parenth{V_{\omega}, \in} &\models \text{ZFC without the Axiom of Infinity}
    \end{align*}
    but the stronger fact that
    \begin{align*}
        \parenth{V_{\omega}, \in} &\models \text{ZFC without the Axiom of Infinity but with its \textit{negation}}
    \end{align*}
    Thus,
    \begin{align*}
        \ZFC \vdash \text{ZFC without Infinity but with $\neg\text{Infinity}$ is consistent}
    \end{align*}
    Ie, from ZFC, we can deduce that the entirety of ZFC, except with the negation of the axiom of infinity instead of the axiom itself, is consistent.
\end{remark}

Here, we end our discussion on the $V$-hierarchy, and proceed to consider a broader type of object than sets, known as classes.
